\section{Results and Discussion} \label{sec:results}

In this work, we investigate the interactions between population size, adaptive potential, and population structure in influencing selection on mutator traits within asexual populations.
In our first experiments, we replicate previous work by \citet{raynes2018sign} showing sign-change effects of population size on mutator favorability.
Subsequent experiments investigate how these dynamics vary based on initial mutator abundance and the availability of beneficial mutations.
Finally, experiments compare mutator dynamics between well-mixed population structure and spatial segregation into subpopulations connected by migration.
[Sentence: throughout, we compare simulation outcomes against the analytic approximation developed in Section \ref{sec:model}, drawn as dashed black lines in figures.]

\subsection{Scaling Population Size}
\label{sec:scaling-population-size}

In a first set of experiments, we tested the capability of our model to reproduce the qualitative regimes of mutator dynamics across population scales identified by \citet{raynes2018sign}.
In the first regime, where population size is very small, fixation of mutator and nonmutator alleles is nearly equiprobable, owing to the overwhelming influence of stochastic effects.
Subsequently, in medium-sized populations, mutators become disfavored on account of increased power of selection becoming in penalizing the increased mutational load imposed by mutator allelees.
Finally, in sufficiently large populations, a sign-change fitness effect occurs and mutator fixation becomes overwhelmingly favored.
This effect occurs in asexual populations because population size increases the probability of at least one mutator discovering a beneficial mutation, and so ``hitch-hiking'' to sweep out nonmutators.

\begin{figure}[h]
  % adapted from https://tex.stackexchange.com/a/122813/316176
  \captionsetup[subfigure]{justification=raggedright}
  \begin{minipage}{0.7\textwidth}

    \begin{minipage}{0.04\textwidth}~\end{minipage}%
    \begin{minipage}{0.44\textwidth}
      \centering
      \itshape
      \textbf{unlimited} beneficial mutations
    \end{minipage}%
    \begin{minipage}{0.34\textwidth}
      \centering
      \itshape
      \textbf{single} beneficial mutation
    \end{minipage}

    ~\vspace{-1.5ex}

    % Top subfigure
    \begin{subfigure}[b]{\linewidth}
      \begin{minipage}{0.88\textwidth}
        \begin{minipage}{0.53\textwidth}
          \includegraphics[height=2.9cm, trim={0cm 0.2cm 3.5cm 0.8cm}, clip]{bindle/bindle-wse-5050-spatial2d-32atile-infben-traits.py/bindle/teeplots/wse-5050-spatial2d-32atile-infben-traits/errorbar=ci+hue=genotype+num-abm=(inf,)+style=genotype+viz=size-fixation-areaplot+x=population-size+y=fixation-probability+ext=.pdf}%
        \end{minipage}%
        \begin{minipage}{0.47\textwidth}
          \includegraphics[height=2.9cm, trim={1.3cm 0.2cm 3.5cm 0.8cm}, clip]{bindle/bindle-wse-5050-spatial2d-32atile-infben-traits.py/bindle/teeplots/wse-5050-spatial2d-32atile-infben-traits/errorbar=ci+hue=genotype+num-abm=(1.0,)+style=genotype+viz=size-fixation-areaplot+x=population-size+y=fixation-probability+ext=.pdf}
        \end{minipage}
      \end{minipage}%
      \hspace{-3ex}%
      \begin{minipage}{0.12\textwidth}
        \raggedright
        \caption{\footnotesize 32 agents per PE\\~\\~\\}
        \label{fig:wse-inf-one:32}
      \end{minipage}%
    \end{subfigure}%

    % Bottom subfigure - Adjusted layout to match the top
    \begin{subfigure}[b]{\linewidth}
      \begin{minipage}{0.88\textwidth}
        \begin{minipage}{0.53\textwidth}
          \includegraphics[height=2.9cm, trim={0cm 0.2cm 3.7cm 0.8cm}, clip]{bindle/bindle-wse-5050-spatial2d-2048atile-infben-traits.py/bindle/teeplots/wse-5050-spatial2d-2048atile-infben-traits/errorbar=ci+hue=genotype+num-abm=(inf,)+style=genotype+viz=size-fixation-areaplot+x=population-size+y=fixation-probability+ext=.pdf}%
        \end{minipage}%
        \begin{minipage}{0.47\textwidth}
          \includegraphics[height=2.9cm, trim={1.3cm 0.2cm 3.7cm 0.8cm}, clip]{bindle/bindle-wse-5050-spatial2d-2048atile-infben-traits.py/bindle/teeplots/wse-5050-spatial2d-2048atile-infben-traits/errorbar=ci+hue=genotype+num-abm=(1.0,)+style=genotype+viz=size-fixation-areaplot+x=population-size+y=fixation-probability+ext=.pdf}
        \end{minipage}
      \end{minipage}%
      \hspace{-3ex}%
      \begin{minipage}{0.1\textwidth}
        \caption{\footnotesize 2,048 agents per PE\\~\\~\\}
        \label{fig:wse-inf-one:2048}
      \end{minipage}%
    \end{subfigure}%

~\vspace{-1.2ex}

\includegraphics[width=0.9\textwidth, trim={0cm 4.7cm 6.2cm 0cm}, clip]{bindle/bindle-wse-5050-spatial2d-2048atile-infben-traits.py/bindle/teeplots/wse-5050-spatial2d-2048atile-infben-traits/col=available-beneficial-mutations+errorbar=sd+hue=genotype+kind=line+style=genotype+viz=relplot+x=population-size+y=fixation-probability+ext=}

  \end{minipage}%
  \begin{minipage}{0.3\textwidth}
    \caption{%
      \textbf{Restricted adaptive potential favors nonmutators in large populations.}
      \footnotesize
      Area plots compare mutator versus nonmutator fixation probabilities across surveyed population sizes.
      When adaptive potential is unlimited (left column), mutators are strongly favored in large populations, owing to their capacity to more rapidly discover beneficial mutations.
      Right column shows fixation outcomes across population sizes, with adaptive potential restricted to just one beneficial mutation.
      As before, mutators gain favor in intermediate population sizes.
      However, under the adaptation-restricted regime, nonmutators regain favor at very large population sizes.
      Top panel, \ref{fig:wse-inf-one:32}, shows results with subpopulation size of 32 agents per PE, scaling population size up to 23.9 million.
      Bottom panel, \ref{fig:wse-inf-one:2048}, reports 2,048 agents per PE with population sizes up to 1.5 billion.
      Experiments were conducted on WSE with counter-based genome model, using populations initialized with a 50/50 mix of non- and mutators.
      Error bands indicate bootstrapped 95\% confidence intervals.
      Supplementary Figure \ref{fig:fixheat-wse-altatile} details results in a tabular format.
      [Dashed black lines show the analytic approximation of Section \ref{sec:model}.]
    }
    \label{fig:wse-inf-one}
  \end{minipage}
\end{figure}


As shown in the left panel of Figure \ref{fig:wse-inf-one:32}, our simulations reproduce the sign-change fitness effect described by \citet{raynes2018sign}.
Leveraging the capabilities of WSE accelerator hardware, we tested population sizes several orders of magnitude beyond those reported in \citet{raynes2018sign}.
For these experiments, 2D population structure was used with 50/50 initialization.
In line with expectations, no further qualitative differences in mutator fixation probability emerged in these larger population sizes, which ranged up to 1 billion agents.
For the largest-scale trials, shown in the left panel of Figure \ref{fig:wse-inf-one:2048}, an increased subpopulation size of 2,048 agents per deme was used.
As such, in these trials, note that coarsened spacing between surveyed population sizes obscures the sign-change fitness effect near the lower bound of population size.

[Paragraph comparing against \cref{eq:escape}: with unlimited adaptive potential the analytic approximation reduces to the escape term, which reproduces the sign change at $N_\mathrm{crit}^{-} = \delta / u_m \approx 600$ and saturation of mutator fixation probability at larger sizes (dashed lines, left panels of Figure \ref{fig:wse-inf-one}); note the lower sign change is unaffected by deme size.]

\subsection{Restricting Adaptive Potential}
\label{sec:restricting-adaptive-potential}

Next, we sought to establish how adaptive potential interacts with current synthesis of the role of population size in mediating mutator dynamics.
One assumption of existing work by \citet{raynes2018sign} is availability of abundant potential beneficial mutations.
Existing simulation studies suggest that restricting adaptive potential should disfavor mutator alleles \citep{tenaillon1999mutators}, but understanding of how this effect interacts with other factors is limited.

To investigate this question, we next performed experiments where the supply of possible beneficial mutations was constrained.
As for the above results, experiments were conducted on the WSE platform using 2D population structure with 50/50 initialization.
The right panels of \cref{fig:wse-inf-one:32,fig:wse-inf-one:2048} show mutator fixation probabilities where agents could accrue no more than one beneficial mutation for trials using 32 and 2,048 agents per deme, respectively.
In both cases, a second sign-change fitness effect for mutator traits appears, with nonmutators regaining favor past population sizes of 100,000 agents.
Beyond population sizes of 1 million agents, nonmutators reliably drive mutators to extinction.

These results substantiate an earlier verbal hypothesis by \citet{desai2011balance} that more than one beneficial mutation should become necessary to fix a mutator allele when fixation time becomes long, as is the case with large population size.
The rationale behind this expectation is that where fixation time is long, non-mutators are more likely to discover a comparable beneficial allele before being driven to extinction.
From this point, disadvantage in mutational load dooms the mutator allele.

[Paragraph: \cref{eq:catchup} makes this argument quantitative for $K = 1$. Report predicted upper critical population size $N_\mathrm{crit}^{+}(1)$ from \cref{eq:ncrit-upper} ($1.1 \times 10^5$ well-mixed; somewhat lower with deme structure) against the observed second sign change near 100,000 agents (dashed lines, right panels of Figure \ref{fig:wse-inf-one}). Note where the approximation is rough: it overshoots peak mutator fixation probability near $N \approx 5{,}000$ and falls off somewhat more sharply than simulation above $N_\mathrm{crit}^{+}$.]

\begin{figure*}

\begin{minipage}{0.65\textwidth}
  \includegraphics[width=\textwidth]{bindle/bindle-wse-5050-spatial2d-2048atile-traits.py/bindle/teeplots/wse-5050-spatial2d-2048atile-traits/col=available-beneficial-mutations+errorbar=ci+hue=genotype+style=genotype+viz=size-fixation-areaplot+x=population-size+y=fixation-probability+ext=.pdf}%
  \end{minipage}
\begin{minipage}{0.3\textwidth}
\caption{
\textbf{Availability of adaptive potential shifts favor from nonmutators to mutators.}
\footnotesize
Area plots show non- vs. mutator fixation probabilities.
As available beneficial mutations are increased, mutators gain favor in progressively larger population sizes.
Simulations were conducted on WSE using the counter-based genome model, with populations initialized to a 50/50 mix of non- and mutators.
Subpopulations comprised 2,048 agents per PE.
Shaded bands show bootstrapped 95\% confidence intervals.
Supplementary Figure \ref{fig:fixheat-wse-altatile:2048} details results in a tabular format.
[Dashed black lines show the analytic approximation of Section \ref{sec:model}.]
}
\label{fig:avail-ben-muts}
\end{minipage}

\end{figure*}


Next, we tested sensitivity of the relationship between mutator fixation probability and population size with respect to the amount of adaptive potential available.
Figure \ref{fig:avail-ben-muts} shows fixation curves where one, two, three, four, and five beneficial mutations are available, conducted with 256 agents per deme.
Under surveyed conditions, with 2D spatial structure and mutators initially in equal proportion to nonmutators, nonmutator persistence rapidly decays with increases in adaptive potential.
Even with the largest surveyed billion-agent population size, mutators regain selective parity where four or more beneficial mutations are available.
Interestingly, this result may help explain an earlier observation by \citet{tenaillon1999mutators} that below a threshold of four beneficial mutations available, mutator alleles never rose from background levels to fix within their compartment-based model of a well-mixed billion-member population.
Notwithstanding differences in model parameterization and population structure, our findings suggest that this threshold may coincide with the point where the net fitness effect of the mutator allele neutralizes in commensurately scaled populations.

[Paragraph comparing Figure \ref{fig:avail-ben-muts} against \cref{eq:fix,eq:ncrit-upper}: the share law predicts $N_\mathrm{crit}^{+}$ rises by about $u_m / u_n = M q_0 \approx 90$ per available beneficial mutation ($N_\mathrm{crit}^{+} \approx 1.1 \times 10^5$, $5.0 \times 10^6$, $4.4 \times 10^8$, $4.0 \times 10^{10}$ for $K = 1$ to 4, well-mixed). Describe agreement (dashed lines) and where the approximation shifts $N_\mathrm{crit}^{+}$ too little under 2D structure at $K \geq 3$. Note the prediction places a billion-member population between $N_\mathrm{crit}^{+}(3)$ and $N_\mathrm{crit}^{+}(4)$, consistent with the four-mutation threshold of \citet{tenaillon1999mutators}.]

\begin{figure*}

\begin{minipage}{0.65\textwidth}
  \includegraphics[width=\textwidth]{bindle/bindle-cupy-5050-traits.py/bindle/teeplots/cupy-5050-traits/col=available-beneficial-mutations+errorbar=ci+hue=population-structure+kind=line+palette=dark2+style=population-structure+viz=relplot+x=population-size+y=fixation-probability+ext=.pdf}%
\end{minipage}
\begin{minipage}{0.3\textwidth}
\caption{
\textbf{Well-mixed and spatially-structured populations exhibit similar relationship between population size and mutator fixation probability in 50/50 competition experiments.}
\footnotesize
In both scenarios, as available beneficial mutations are increased, mutators gain favor in progressively larger population sizes.
Simulations were conducted on GPU using the counter-based genome model, with populations initialized to a 50/50 mix of non- and mutators.
Subpopulations comprised 256 agents per PE.
Shaded bands show bootstrapped 95\% confidence intervals.
[Dashed black lines show the analytic approximation of Section \ref{sec:model}.]
}
\label{fig:avail-ben-muts-gen}
\end{minipage}

\end{figure*}


Given that experiments conducted on the WSE platform all involve 2D spatial structure, we implemented an additional set of GPU-based experiments to compare findings with well-mixed population structure.
Figure \ref{fig:avail-ben-muts-gen} overviews these results.
Although experiments ranged up to a smaller maximum population size on the GPU platform, detected regimes of mutator dynamics are consistent with WSE-based experiments.
Comparing on the basis of spatial structure, we find fixation probabilities to be near-identical between GPU-based trials using 2D and well-mixed structure across surveyed conditions.
% https://github.com/mmore500/hypermutator-dynamics/blob/484ac79f35baf10e4acc52babee836d10ab50bdb/binder/cupy-5050-traits.ipynb
We find significant divergence in fixation probability between spatial structure treatments only for population size 20,736 with adaptive potential of one available beneficial mutation (Mann-Whitney U test; $p < 0.0005$; Bonferroni corrected $p < 0.01$).

[Paragraph: the analytic approximation accounts for this near-identity. Under 50/50 initialization, structure enters only through sweep time (\cref{eq:sweep-mixed,eq:sweep-2d}), which shifts predicted $N_\mathrm{crit}^{+}$ by less than a factor of two for 256-agent demes (dashed lines in Figure \ref{fig:avail-ben-muts-gen}, one per structure).]

\subsection{Effects of Background Mutator Prevalence}
\label{sec:background-hypermutator-prevalence}

\begin{figure}[h]
\begin{minipage}{0.6\textwidth}
  \includegraphics[width=\linewidth]{bindle/bindle-wse-traits.py/bindle/teeplots/wse-traits/errorbar=ci+hue=initial-conditions+kind=line+palette=set2-r+row=population-size+style=initial-conditions+viz=relplot+what=counter-based+x=available-beneficial-mutations+y=fix-prob+ext=.pdf}%
\end{minipage}%
\begin{minipage}{0.4\textwidth}
  \caption{
  \textbf{Initial supply of mutators influences nonmutator resilience to available adaptive potential.}
  \footnotesize
  Lineplot strips show change in mutator fixation probability with adaptive potential for surveyed population sizes, with error bands providing bootstrapped 95\% confidence intervals.
  Simulations were conducted on WSE, using a population size of 256 agents per PE and counter-based genome model.
  Results using a site-explicit model were similar, provided in Supplementary \cref{fig:denovo-5050-conditions-combined-site-explicit,fig:wse-site-explicit-counter-based}.
  Supplementary Figure \ref{fig:fixheat-wse-256atile} details results in a tabular format.
  [Dashed black lines show the analytic approximation of Section \ref{sec:model}.]
    }
    \label{fig:denovo-5050-conditions-combined}
  \end{minipage}
\end{figure}


Having observed large population size prevent mutator fixation in competition experiments under conditions of limited adaptive potential, we next sought to assess how such an effect might manifest in more naturalistic conditions.
For this purpose, we relaxed the initial condition that mutators begin ``50/50'' in equal proportion to nonmutators.
Instead, we allowed mutators to enter the population at a low background rate \citep{desai2011balance,johnson1999approach}, an experimental design common with several previous simulation studies of hypermutator dynamics \citep{wylie2009fixation,tenaillon1999mutators}.
To provide a conservative reflection of the impact of mutator allele rarity, we used a mutator allele introduction probability of $10^{-6}$ per replication event.
Figure \ref{fig:denovo-5050-conditions-combined} compares mutator fixation probabilities between 50/50 and \textit{de novo} conditions.
Trials were conducted with 2D spatial structure on the WSE platform.
With 256 agents per deme, population size scaled up to 136 million agents.

Prior work analyzing conditions with infinite adaptive potential has shown per-capita favorability of mutator traits to be density independent \citep{raynes2019selection}.
As such, a lower background rate for mutator alleles should be expected to substantially reduce the net probability of mutator fixation.
In line with expectations, the \textit{de novo} treatment substantially reduces mutator fixation.
Under \textit{de novo} conditions, nonmutators reliably resist mutator fixation in populations of 136 million through more than 10 beneficial mutations available.
Under 50/50 conditions, by contrast, mutators begin to fix in these large populations past 2 beneficial mutations available.

[Paragraph comparing Figure \ref{fig:denovo-5050-conditions-combined} against Section \ref{sec:model:denovo}: in small populations, fixation probability follows $1 - (1 - P_\mathrm{sweep})^K$ with $P_\mathrm{sweep} \approx 0.009$ independent of population size (\cref{eq:denovo-small}); in large populations it switches near the critical adaptive potential $K^*$ of \cref{eq:kstar-2d}, predicted at 12 to 13.5 available beneficial mutations against 12 to 16 observed. Note the approximation places the switch about one mutation early at the largest sizes.]

We performed sensitivity analysis of our findings to the underlying counter-based genome model used.
For this purpose, we implemented an alternate ``site-explicit'' model where the probability of discovering a beneficial mutation scales proportional to the number available.
As shown in Supplementary Figure \ref{fig:wse-site-explicit-counter-based}, results were generally consistent between the two genome models.
Notably, though, under the ``site-explicit'' model, mutator alleles are more frequently fixed within the smallest surveyed population size under 50/50 initialization.

\subsection{Spatiotemporal Dynamics}
\label{sec:dynamics}

To assess how competition between non- and mutator strains unfolds at very large population scales, we performed an additional set of simulations on the WSE using the site-explicit model of adaptation, where mutators originated \textit{de novo} from nonmutator strains.
For these experiments, a subpopulation size of 256 agents per PE was used to preserve on-device memory to record time series data.
This configuration yielded a net population size of 191 million agents.
In these trials, we probed the critical region where outcomes transition between nonmutator and mutator fixation, configuring adaptive potential between 12 and 16 available beneficial mutations.

\begin{figure}[ht]
  \centering
  \resizebox{\columnwidth}{!}{%
  \begin{minipage}{\textwidth} %

  %----------------------------
  % First row: NO FIXATION
  \begin{subfigure}[t]{\textwidth}
    \centering
    % Left minipage for row label (empty caption)
    \begin{minipage}[b]{0.04\textwidth}
      \caption{}
      \label{fig:dynamics:no-fixation}
    \end{minipage}
    % Six images (columns)
    \begin{minipage}[b]{0.15\textwidth}
      \includegraphics[width=\linewidth]{binder/binder-wse-denovo-spatial2d-explicitsite-timeseries.ipynb/binder/wse-denovo-spatial2d-explicitsite-timeseries/a=traitframes+nmut=14+rep=39a89ca6-a1b5-4b32-ae5f-f0dbb40ba027/dstream_Tbar=000504+ext=.png}
    \end{minipage}
    \begin{minipage}[b]{0.15\textwidth}
      \includegraphics[width=\linewidth]{binder/binder-wse-denovo-spatial2d-explicitsite-timeseries.ipynb/binder/wse-denovo-spatial2d-explicitsite-timeseries/a=traitframes+nmut=14+rep=5dc8e084-0382-4d7b-9b76-6c3902ca3c1d/dstream_Tbar=000952+ext=.png}
    \end{minipage}
    \begin{minipage}[b]{0.15\textwidth}
      \includegraphics[width=\linewidth]{binder/binder-wse-denovo-spatial2d-explicitsite-timeseries.ipynb/binder/wse-denovo-spatial2d-explicitsite-timeseries/a=traitframes+nmut=14+rep=5dc8e084-0382-4d7b-9b76-6c3902ca3c1d/dstream_Tbar=001400+ext=.png}
    \end{minipage}
    \begin{minipage}[b]{0.15\textwidth}
      \includegraphics[width=\linewidth]{binder/binder-wse-denovo-spatial2d-explicitsite-timeseries.ipynb/binder/wse-denovo-spatial2d-explicitsite-timeseries/a=traitframes+nmut=14+rep=5dc8e084-0382-4d7b-9b76-6c3902ca3c1d/dstream_Tbar=002552+ext=.png}
    \end{minipage}
    \begin{minipage}[b]{0.15\textwidth}
      \includegraphics[width=\linewidth]{binder/binder-wse-denovo-spatial2d-explicitsite-timeseries.ipynb/binder/wse-denovo-spatial2d-explicitsite-timeseries/a=traitframes+nmut=14+rep=5dc8e084-0382-4d7b-9b76-6c3902ca3c1d/dstream_Tbar=020472+ext=.png}
    \end{minipage}
    \begin{minipage}[b]{0.15\textwidth}
      \includegraphics[width=\linewidth]{binder/binder-wse-denovo-spatial2d-explicitsite-timeseries.ipynb/binder/wse-denovo-spatial2d-explicitsite-timeseries/a=traitframes+nmut=14+rep=5dc8e084-0382-4d7b-9b76-6c3902ca3c1d/dstream_Tbar=061424+ext=.png}
    \end{minipage}
  \end{subfigure}

  \vspace{1em} % vertical space between rows

  %----------------------------
  % Second row: FIXATION
  \begin{subfigure}[t]{\textwidth}
    \centering
    % Left minipage for row label (empty caption)
    \begin{minipage}[b]{0.05\textwidth}
      \caption{}
      \label{fig:dynamics:fixation}
    \end{minipage}
    % Six images (columns)
    \begin{minipage}[b]{0.15\textwidth}
      \includegraphics[width=\linewidth]{binder/binder-wse-denovo-spatial2d-explicitsite-timeseries.ipynb/binder/wse-denovo-spatial2d-explicitsite-timeseries/a=traitframes+nmut=14+rep=39a89ca6-a1b5-4b32-ae5f-f0dbb40ba027/dstream_Tbar=000504+ext=.png}
    \end{minipage}
    \begin{minipage}[b]{0.15\textwidth}
      \includegraphics[width=\linewidth]{binder/binder-wse-denovo-spatial2d-explicitsite-timeseries.ipynb/binder/wse-denovo-spatial2d-explicitsite-timeseries/a=traitframes+nmut=14+rep=39a89ca6-a1b5-4b32-ae5f-f0dbb40ba027/dstream_Tbar=000952+ext=.png}
    \end{minipage}
    \begin{minipage}[b]{0.15\textwidth}
      \includegraphics[width=\linewidth]{binder/binder-wse-denovo-spatial2d-explicitsite-timeseries.ipynb/binder/wse-denovo-spatial2d-explicitsite-timeseries/a=traitframes+nmut=14+rep=39a89ca6-a1b5-4b32-ae5f-f0dbb40ba027/dstream_Tbar=001400+ext=.png}
    \end{minipage}
    \begin{minipage}[b]{0.15\textwidth}
      \includegraphics[width=\linewidth]{binder/binder-wse-denovo-spatial2d-explicitsite-timeseries.ipynb/binder/wse-denovo-spatial2d-explicitsite-timeseries/a=traitframes+nmut=14+rep=39a89ca6-a1b5-4b32-ae5f-f0dbb40ba027/dstream_Tbar=002489+ext=.png}
    \end{minipage}
    \begin{minipage}[b]{0.15\textwidth}
      \includegraphics[width=\linewidth]{binder/binder-wse-denovo-spatial2d-explicitsite-timeseries.ipynb/binder/wse-denovo-spatial2d-explicitsite-timeseries/a=traitframes+nmut=14+rep=39a89ca6-a1b5-4b32-ae5f-f0dbb40ba027/dstream_Tbar=002489+ext=.png}
    \end{minipage}
    \begin{minipage}[b]{0.15\textwidth}
      \includegraphics[width=\linewidth]{binder/binder-wse-denovo-spatial2d-explicitsite-timeseries.ipynb/binder/wse-denovo-spatial2d-explicitsite-timeseries/a=traitframes+nmut=14+rep=39a89ca6-a1b5-4b32-ae5f-f0dbb40ba027/dstream_Tbar=002489+ext=.png}
    \end{minipage}
  \end{subfigure}

  %----------------------------
  % Timepoint labels below the bottom row:
  \begin{minipage}[c]{0.154\textwidth}
\hfill
\begin{varwidth}{\textwidth}
$T = 504$
\end{varwidth}
\hfill
  \end{minipage}
  \begin{minipage}[c]{0.154\textwidth}
\hfill
\begin{varwidth}{\textwidth}
$T = 952$
\end{varwidth}
\hfill
  \end{minipage}
  \begin{minipage}[c]{0.154\textwidth}
\hfill
\begin{varwidth}{\textwidth}
$T = 1,400$
\end{varwidth}
\hfill
  \end{minipage}
  \begin{minipage}[c]{0.154\textwidth}
\hfill
\begin{varwidth}{\textwidth}
$T = 2,552$
\end{varwidth}
\hfill
  \end{minipage}
  \begin{minipage}[c]{0.154\textwidth}
\hfill
\begin{varwidth}{\textwidth}
$T = 20,472$
\end{varwidth}
\hfill
  \end{minipage}
  \begin{minipage}[c]{0.154\textwidth}
\hfill
\begin{varwidth}{\textwidth}
$T = 61,424$
\end{varwidth}
\hfill
  \end{minipage}

  \end{minipage}%
  }%

  \caption{
  \textbf{Spatiotemporal composition of simulated populations from Wafer-Scale Engine experiments.}
  \footnotesize
  Snapshots show 191 million agent populations with 2D spatial structure, using site-explicit genome model configured to adaptive potential of 14 beneficial mutations available and mutators introduced \textit{de novo}.
  Raster values are binary, with white pixels indicating a sampled nonmutators and black pixels indicating a sampled mutator.
  Subpanels \ref{fig:dynamics:no-fixation} and       \ref{fig:dynamics:fixation} show replicates where mutators do not, and do, reach fixation, respectively.
  Animations are provided at \url{https://hopth.ru/ej} and \url{https://hopth.ru/ek}.
  Example timecourses for simulations with 12 and 16 beneficial mutations available can be found at \url{https://hopth.ru/el} and \url{https://hopth.ru/em}.
  }
  \label{fig:dynamics}
\end{figure}


Figure \ref{fig:dynamics} arranges sequential population snapshots from two example simulations, both with adaptive potential of 14 beneficial mutations.
In the first example, mutators fixed; in the second, they did not.
At outset, panels \ref{fig:dynamics:no-fixation} and \ref{fig:dynamics:fixation} unfold similarly.
First, mutator strains appear \textit{de novo} across the breadth of simulated populations.
Then, as patches of mutators expand outwards, the nonmutator population is broken into pockets.
Mutator strains continued to stem from nonmutator lineages as they accrued adaptive mutations.
Of 16 replicates where we observed mutators fix, we found 2 cases where the final dominant mutator lineage arose from a partially adapted nonmutator lineage --- harboring 1 and 3 beneficial mutations, respectively.
% https://github.com/mmore500/hypermutator-dynamics/blob/90422b30f2dd9ca29baad273cbcc676dbfeab55f/binder/wse-denovo-spatial2d-explicitsite-genomes.ipynb
In scenarios where mutators fix, surviving patches of nonmutators dwindle away entirely (panel \ref{fig:dynamics:fixation}).
Where nonmutators persist, they re-establish via concentric growth from surviving pockets (panel \ref{fig:dynamics:no-fixation}).

\begin{figure}

\begin{minipage}{0.4\textwidth}
\includegraphics[width=\textwidth]{bindle/bindle-wse-denovo-spatial2d-explicitsite-timeseries-2025-02-07-peaksweep.py/bindle/teeplots/wse-denovo-spatial2d-explicitsite-timeseries-2025-02-07-peaksweep/color=orange+kind=line+viz=relplot+x=cerebraslib-hypermut-num-avail-ben-muts+y=dstream-value+ext=.pdf}
\end{minipage}
\begin{minipage}{0.55\textwidth}
\caption{
\textbf{Mutators prevalence transitively peaks above population majority in scenarios with limited adaptive potential.}
\footnotesize
Simulations were conducted on WSE, configured with population size 191 million over 2D spatial structure and site-explicit genome model with mutators introduced \textit{de novo}.
Under these conditions, mutator fixation was observed only for 14 or more available beneficial mutations.
Shaded bands show 100\% percentile interval across four replicates.
[Dashed black lines show the analytic approximation of Section \ref{sec:model}.]
}
\label{fig:peaksweep}
\end{minipage}
\end{figure}


At peak mutator prevalence, we observed nonmutators shrink to less than 0.1\% of the population in experiments with adaptive potential of 14 available beneficial mutations.
% ^^^ https://github.com/mmore500/hypermutator-dynamics/blob/5e904ac2fbbd4b41a0ad679883fc1a63af71c00c/binder/wse-denovo-spatial2d-explicitsite-timeseries.ipynb
To assess the magnitude of transient peak mutator concentrations across levels of adaptive potential, we performed an additional set of four-replicate simulations with adaptive potentials ranging from 2 to 10 available beneficial mutations.
For all replicates with 6 or more available beneficial mutations, mutators peaked beyond 50\% prevalence (Figure \ref{fig:peaksweep}).

[Paragraph: nucleation-and-growth interpretation of Figure \ref{fig:dynamics} (mutator patches founded \textit{de novo} and expanding at a wave speed, racing nonmutator adaptation) and comparison of observed peak prevalence against \cref{eq:peak-prevalence} (dashed line in Figure \ref{fig:peaksweep}; predicted 0.04, 0.28, 0.69, 0.94, 1.00 at 2 to 10 available beneficial mutations). Caveat that these simulations used the site-explicit genome model whereas the approximation assumes the counter-based one.]

\subsection{Effects of Population Structure}
\label{sec:population-structure}

In a final set of experiments, we extended our earlier trials with 50/50 initialization to assess the impact of spatial structure on mutator dynamics where mutators are initially rare.
In particular, we were curious whether removal of population structure would negligibly affect mutator fixation probabilities, as was earlier the case with 50/50 initialization.
Trials conducted in this experiment comparing the 2D deme structure to fully well-mixed conditions were conducted using GPU evaluation.
As a result, population sizes ranged only up to 15 million.

\begin{figure}[h]
\begin{minipage}{\textwidth}
  \includegraphics[width=\linewidth]{bindle/bindle-cupy-traits.py/bindle/teeplots/cupy-traits/errorbar=ci+exclude=1D-demes+hue=population-structure+kind=line+palette=dark2+row=population-size+style=initial-conditions+viz=relplot+x=available-beneficial-mutations+y=fix-prob+ext=.pdf}%
\end{minipage}

\begin{minipage}{\textwidth}
  \caption{%
    \textbf{Population structure increases nonmutator resilience to available adaptive potential.}
    \footnotesize
    Lineplot strips show relationship between adaptive potential and mutator fixation probability across surveyed population sizes, with error bands indicating bootstrapped 95\% confidence intervals.
    In the 50/50 treatment, experiments were initialized with an even mix of mutator and nonmutator agents.
    In the \textit{de novo} treatment, mutators are initially absent and arise spontaneously \textit{de novo} with probability $10^{-6}$.
    Effects of population structure are more pronounced under \textit{de novo} conditions, when initial supply of mutators is scarce.
    Under well-mixed conditions, increase in population size favors mutators; nonmutators reliably persist in very large populations through fewer than 10 available beneficial mutations.
    In contrast, under 2D deme structure, very large population sizes are favorable to nonmutators.
    With such strong spatial structure, very large populations of nonmutators resist mutator fixation even when upwards of 10 beneficial mutations are available.
    Results are under GPU simulation, using counter-based genome model.
    Supplementary \cref{fig:fixheat-5050-cupy,fig:fixheat-denovo-cupy} detail results in a tabular format.
    [Dashed black lines show the analytic approximation of Section \ref{sec:model}.]
  }
  \label{fig:spatial-structure-combined}
\end{minipage}
\end{figure}


Figure \ref{fig:spatial-structure-combined} plots increase in mutator fixation probability with available adaptive potential across a spectrum of population sizes.
Interestingly, spatial structure does cause results to differ substantially under \textit{de novo} conditions.
We find that in large nonmutator populations of 15 million with well-mixed structure, fixation probability of mutator alleles becomes appreciable with adaptive potential of only 5 available beneficial mutations.
Thisis not the case for large populations with 2D deme structure, where fixation of mutator alleles become appreciable above 10 available beneficial mutations.
For intermediate population sizes of 186,000 agents, however, an opposite effect occurs: spatial structure raises the probability of mutator fixation across the board of available adaptive potential compared to well-mixed conditions.

These findings contrast with previous modeling work assuming unlimited adaptive potential, which predicts that so long as any level of connectivity exists between subpopulations, spatial structure boosts mutator fixation probability \citep{raynes2019migration}.
One possible explanation for why spatial structure can decrease mutator fixation probability is in extending fixation times, which extends the opportunity for nonmutators to exploit available adaptive potential before being driven to extinction.
This described effect notably aligns with findings from \textit{in vivo} evolution experiments with \textit{E. coli} that show that migration barriers delaying introduction of mutator strains into nonmutator populations can reduce mutator advantage \citep{lechat2006escherichia}.

[Paragraph: the analytic approximation supports this explanation and makes it quantitative (dashed lines in Figure \ref{fig:spatial-structure-combined}). Well-mixed, $K^*$ follows \cref{eq:kstar-mixed} (about 4.5 to 5.5, against 6 to 7 in simulation); with 2D structure, \cref{eq:kstar-2d} (about 12 to 13.5). The opposite effect at intermediate sizes also follows: on small grids a single mutator domain covers the lattice before nonmutators adapt, while slower sweeps give each mutator more time to hitchhike.]

On a separate front, it is surprising that substantial differences in mutator dynamics between deme-structured and well-mixed conditions only become apparent with low initial mutator frequency, given the established expectation of fitness effects of mutator alleles as a density- and number- independent \citep{raynes2019selection}.
While one possible explanation may be in part that detectability of fitness effect magnitudes is masked by sensitivity limitations of the 50/50 assay, this result draws into question the extent to which density-independence of mutator fitness effects generalizes to scenarios with limited adaptive potential and spatial structure.
In future work, it would be informative to test whether the impact of spatial structure where mutators are rare stems from heterogeneity in the spatial distribution of mutator clades or from slowing the spread of beneficial alleles.
[Sentence or two: in the analytic model the difference arises because, under 50/50 initialization, mutators are present everywhere and structure only modestly lengthens sweeps, whereas mutators arising \textit{de novo} must additionally spread spatially from where they are founded; i.e., the effect stems from spatial heterogeneity of mutator clades coupled with slowed spread.]

\subsection{[Application to Experimental Data]}
\label{sec:callens}

[Optional subsection applying Section \ref{sec:model:denovo} to the evolution experiments of \citet{callens2023hypermutator}: with the model clock calibrated from sweeps completed in nonmutator populations and a lag for detection, a single mutator supply rate of about $10^{-7}$ per generation shared across treatments reproduces hypermutator emergence across control, antibiotic (23 adaptive loci) and osmotic (43 adaptive loci) treatments; a likelihood ratio test finds no evidence for treatment-specific supply rates. Caveats: one fitted parameter, and sensitivity to assumed deleterious mutation rate and mutator strength. See Figure \ref{fig:callens}.]

\begin{figure}[h]
\begin{minipage}{0.5\textwidth}
  \includegraphics[width=\linewidth]{bindle/bindle-callens2023-hypermutator-emergence.py/bindle/teeplots/callens2023-hypermutator-emergence/viz=plot-takeover-by-stress+ext=.pdf}%
\end{minipage}%
\begin{minipage}{0.5\textwidth}
  \caption{
  TODO
  }
  \label{fig:callens}
\end{minipage}
\end{figure}

